\subsection{Thiết kế Data Ingest và Data Preprocessing}

Data Ingest Pipeline là lớp đầu vào của RAG, chịu trách nhiệm tiếp nhận và xử lý tài liệu. Trong hệ thống tuyển sinh, tài liệu đầu vào có thể là PDF, DOCX, XLSX hoặc các biểu mẫu hành chính. Đây là nhóm dữ liệu không đồng nhất, có thể chứa văn bản, bảng, tiêu đề, danh sách, mã ngành, tổ hợp môn, điểm chuẩn và mốc thời gian.

Trước khi tài liệu được đưa vào kho tri thức, hệ thống phải kiểm tra nguồn dữ liệu, đăng ký ingest, parse nội dung, kiểm tra OCR nếu cần, chuẩn hóa Unicode, kiểm tra metadata và chỉ cho phép các tài liệu đạt chất lượng đi tiếp sang bước chunking và indexing.

\begin{figure}[H]
    \centering
    \includegraphics[width=0.95\textwidth]{\detokenize{content/source/RAG/Pipeline Data Ingest và Data Preprocessing.png}}
    \caption{Pipeline data ingest và preprocessing cho kho tri thức RAG tuyển sinh}
    \label{fig:rag-data-ingest-preprocessing}
\end{figure}

Luồng ingest tài liệu quan sát được gồm các bước: admin upload tài liệu hoặc scan thư mục input, pipeline đọc file PDF/DOCX/XLSX, trích xuất văn bản, chuẩn hóa, chia chunks, sinh embedding bằng BAAI/bge-m3, lưu metadata và trạng thái vào RedisDocStatusStorage, lưu vector vào Qdrant và lưu quan hệ tri thức vào Memgraph.

Mỗi tài liệu cần có metadata tối thiểu:

\begin{itemize}
    \item \texttt{document\_id}, \texttt{file\_name}, \texttt{file\_type}, \texttt{source\_type}.
    \item \texttt{created\_at}, \texttt{updated\_at}, \texttt{status}.
    \item \texttt{chunk\_count}, \texttt{content\_length}, \texttt{checksum}.
    \item \texttt{track\_id}, \texttt{valid\_status}.
\end{itemize}

Mỗi chunk cần có metadata riêng:

\begin{itemize}
    \item \texttt{chunk\_id}, \texttt{document\_id}, \texttt{chunk\_index}.
    \item \texttt{section\_title}, \texttt{page\_number}, \texttt{content}.
    \item \texttt{content\_length}, \texttt{embedding\_model}, \texttt{created\_at}, \texttt{source\_file}.
\end{itemize}

Vai trò của metadata là giúp hệ thống truy vết nguồn, lọc theo năm, loại tài liệu, trạng thái hiệu lực và hỗ trợ citation khi trả lời. Nếu không có metadata, LLM có thể trả lời đúng về mặt ngữ nghĩa nhưng không chứng minh được câu trả lời lấy từ đâu.
